%% ============================================================
%% Font to Binary Converter プログラム設計仕様書
%% ============================================================
%% ビルド: latexmk main.tex (.latexmkrc で uplatex + dvipdfmx)
%% ============================================================
\documentclass[a4paper,11pt,dvipdfmx]{jsarticle}

%% ─── パッケージ ───────────────────────────────
\usepackage[dvipdfmx]{graphicx}
\usepackage[dvipdfmx]{xcolor}
\usepackage[dvipdfmx,
  bookmarks=true,
  bookmarksnumbered=true,
  colorlinks=true,
  linkcolor={blue!70!black},
  citecolor={green!50!black},
  urlcolor={blue!60!black},
  pdfauthor={Shiozawa},
  pdftitle={Font to Binary Converter プログラム設計仕様書},
  pdfsubject={ドットフォント→バイナリ配列変換Webツール},
]{hyperref}
\usepackage{pxjahyper}
\usepackage[top=25mm,bottom=25mm,left=20mm,right=20mm]{geometry}
\usepackage{fancyhdr}
\usepackage{lastpage}
\usepackage{tabularx}
\usepackage{booktabs}
\usepackage{longtable}
\usepackage{enumitem}
\usepackage{titlesec}
\usepackage{needspace}
\usepackage{tcolorbox}
\tcbuselibrary{skins,breakable}
\usepackage{tikz}
\usetikzlibrary{arrows.meta,positioning,shapes.geometric,fit,calc}
\usepackage{listings}
\usepackage{plistings}
\usepackage{caption}
\usepackage{float}
\usepackage{placeins}
\usepackage{amsmath}

%% ─── 色定義 ──────────────────────────────────
\definecolor{headerblue}{HTML}{1B3A5C}
\definecolor{accentblue}{HTML}{2E86C1}
\definecolor{lightgray}{HTML}{F2F3F4}
\definecolor{codebg}{HTML}{F7F9FB}
\definecolor{warnred}{HTML}{E74C3C}
\definecolor{noteyellow}{HTML}{F39C12}
\definecolor{okgreen}{HTML}{27AE60}

%% ─── 改ページ制御 ─────────────────────────────
\widowpenalty=10000
\clubpenalty=10000
\predisplaypenalty=10000
\interfootnotelinepenalty=10000
\brokenpenalty=5000

%% ─── ヘッダー・フッター ───────────────────────
\pagestyle{fancy}
\fancyhf{}
\fancyhead[L]{\small\color{headerblue}\textbf{Font to Binary Converter} プログラム設計仕様書}
\fancyhead[R]{\small\color{headerblue}Rev.\,1.0}
\fancyfoot[C]{\small\thepage\,/\,\pageref{LastPage}}
\fancyfoot[R]{\small\color{gray}2026-04-21}
\renewcommand{\headrulewidth}{0.8pt}
\renewcommand{\footrulewidth}{0.4pt}

%% ─── セクション見出し（孤立防止付き）─────────
\titleformat{\section}
  {\needspace{5\baselineskip}\Large\bfseries\color{headerblue}}
  {\thesection}{1em}{}
  [\vspace{-0.5em}\textcolor{accentblue}{\rule{\textwidth}{1.2pt}}]
\titleformat{\subsection}
  {\needspace{4\baselineskip}\large\bfseries\color{headerblue}}
  {\thesubsection}{1em}{}
\titleformat{\subsubsection}
  {\needspace{3\baselineskip}\normalsize\bfseries\color{accentblue}}
  {\thesubsubsection}{1em}{}

%% ─── listings 設定 ────────────────────────────
\lstdefinestyle{tsstyle}{
  basicstyle=\ttfamily\small,
  keywordstyle=\color{accentblue}\bfseries,
  commentstyle=\color{gray}\itshape,
  stringstyle=\color{okgreen},
  numberstyle=\tiny\color{gray},
  numbers=left,
  numbersep=8pt,
  frame=single,
  framerule=0.5pt,
  rulecolor=\color{accentblue!40},
  backgroundcolor=\color{codebg},
  breaklines=true,
  breakatwhitespace=false,
  tabsize=2,
  showstringspaces=false,
  captionpos=b,
  xleftmargin=15pt,
  framexleftmargin=15pt,
  morekeywords={const,let,var,function,export,import,from,type,interface,
    return,if,else,for,while,async,await,class,extends,new,this,static,
    public,private,null,undefined,true,false,number,string,boolean,void,
    React,useState,useEffect,useMemo,useReducer,useRef,useCallback},
}
\lstset{style=tsstyle}

%% ─── tcolorbox スタイル ──────────────────────
\tcbset{before={\par\medskip\needspace{2\baselineskip}\noindent}}

\newtcolorbox{apibox}[1]{
  enhanced,
  colback=white,
  colframe=accentblue,
  coltitle=white,
  fonttitle=\bfseries\large,
  title=#1,
  top=4mm,
  arc=2mm,
  breakable,
  before upper={\parindent0pt},
  attach boxed title to top left={yshift=-3mm,xshift=5mm},
  boxed title style={colback=accentblue,arc=1.5mm},
}

\newtcolorbox{cautionbox}[1][注意]{
  enhanced,
  colback=warnred!5,
  colframe=warnred!70,
  coltitle=white,
  fonttitle=\bfseries,
  title={\raisebox{-0.1em}{\textbf{⚠}} #1},
  arc=1.5mm,
  breakable,
  attach boxed title to top left={yshift=-2mm,xshift=5mm},
  boxed title style={colback=warnred!70,arc=1mm},
}

\newtcolorbox{notebox}[1][補足]{
  enhanced,
  colback=accentblue!5,
  colframe=accentblue!50,
  coltitle=white,
  fonttitle=\bfseries,
  title={$\triangleright$ #1},
  arc=1.5mm,
  breakable,
  attach boxed title to top left={yshift=-2mm,xshift=5mm},
  boxed title style={colback=accentblue!50,arc=1mm},
}

\newtcolorbox{funcblock}[1]{
  enhanced,
  colback=white,
  colframe=accentblue!25,
  coltitle=headerblue,
  fonttitle=\bfseries\ttfamily\small,
  title={◆ #1},
  arc=1mm,
  boxrule=0.4pt,
  breakable,
  left=3mm, right=3mm, top=2mm, bottom=2mm,
  before skip=3mm, after skip=3mm,
  before upper={\parindent0pt},
}

\newcommand{\apisection}[1]{%
  \vspace{4mm}%
  {\normalsize\bfseries\color{headerblue} #1}%
  \par\vspace{-1mm}%
  {\color{accentblue!40}\rule{\linewidth}{0.5pt}}%
  \par\vspace{2mm}%
}

%% ─── ドキュメント情報 ─────────────────────────
\newcommand{\doctitle}{Font to Binary Converter\\プログラム設計仕様書}
\newcommand{\docsubtitle}{ドットフォント→バイナリ配列変換 Webツール}
\newcommand{\docauthor}{Shiozawa}
\newcommand{\docversion}{1.0}
\newcommand{\docdate}{2026年4月21日}

%% ============================================================
\begin{document}

%% ─── 表紙 ─────────────────────────────────────
\begin{titlepage}
\thispagestyle{empty}
\begin{center}
\vspace*{30mm}

{\color{accentblue}\rule{\textwidth}{2pt}}
\vspace{8mm}

{\Huge\bfseries\color{headerblue}\doctitle\par}
\vspace{10mm}

{\Large\color{accentblue}\docsubtitle\par}
\vspace{6mm}

{\color{accentblue}\rule{\textwidth}{1pt}}
\vspace{40mm}

\begin{tabular}{r l}
\large\bfseries\color{headerblue} プロジェクト    & \large Font to Binary Converter \\[1mm]
\large\bfseries\color{headerblue} ドキュメント種別 & \large プログラム設計仕様書 \\[1mm]
\large\bfseries\color{headerblue} 版数             & \large Rev.~\docversion \\[1mm]
\large\bfseries\color{headerblue} 発行日           & \large \docdate \\[1mm]
\large\bfseries\color{headerblue} 著者             & \large \docauthor \\
\end{tabular}

\vfill

{\color{gray}\small 本書は Font to Binary Converter の設計仕様を記述する。}
\end{center}
\end{titlepage}

%% ─── 改訂履歴 ────────────────────────────────
\section*{改訂履歴}
\addcontentsline{toc}{section}{改訂履歴}

\begin{tabularx}{\textwidth}{c c c X}
\toprule
\textbf{版} & \textbf{日付} & \textbf{著者} & \textbf{変更内容} \\
\midrule
1.0 & 2026-04-21 & Shiozawa & 初版作成。Python版から Web (React + Vite + TypeScript) 実装へ全面刷新。\\
\bottomrule
\end{tabularx}

\vspace{4mm}

\begin{notebox}[旧版との関係]
本書の Rev.~1.0 は、旧来の Python/Tkinter 実装（\texttt{font\_to\_binary\_GUI.py} および
\texttt{font\_to\_binary\_CUI.py}）を置き換えた Web 実装の仕様である。
旧ソースは \texttt{legacy/} ディレクトリにそのまま残されている。
\end{notebox}

\newpage

\tableofcontents
\newpage

\section{概要}

Font to Binary Converter は、ドットフォントで描画された文字を
二次元バイナリ配列（0/1）または各種エンコード形式に変換し、
C/C++/Python/JavaScript 等のソースコードに直接埋め込める形式で
出力する Web ツールである。

本ツールは以下の用途を主な想定としている。

\begin{itemize}[leftmargin=*,itemsep=2pt]
  \item マイコン（Arduino, ESP32, RZ/A1 等）の 0.96'' OLED、LCD 等
        小型ディスプレイ上へのフォント描画用ビットマップデータ生成。
  \item LED マトリクスパネル（8×8, 16×16, 32×32）の文字表示用データ。
  \item ゲーム・デモ用の固定幅ドットフォント素材生成。
  \item ピクセルアート作成時の下書きテンプレート。
\end{itemize}

\subsection{主な特徴}

\begin{itemize}[leftmargin=*,itemsep=2pt]
  \item \textbf{ブラウザ単体で動作}: サーバー不要。GitHub Pages で静的配信。
  \item \textbf{2モードUI}:「かんたん」は3ステップ(文字/見た目/用途)で
        初心者でも迷わず変換可能。「詳細」は全設定を直接編集可能。
        モード選択は LocalStorage に自動保存。
  \item \textbf{用途別プリセット}: Arduino/C言語, C言語2D, Python リスト,
        MicroPython bytes, JSON, 記号アートの6種を1クリック切替。
  \item \textbf{フォント同梱}: \texttt{DotGothic16} (16×16) と
        \texttt{Misaki Gothic 2nd} (8×8) を同梱。
  \item \textbf{ユーザーフォント対応}: TTF/OTF/WOFF をその場でアップロード可能。
  \item \textbf{多言語出力}: C, C++, Arduino(PROGMEM), Python, JavaScript,
        TypeScript, Rust, Go, JSON, Markdown, プレーンテキスト, 記号アート。
  \item \textbf{詳細な出力設定}: 2D/3D/フラット/ビットパック構造、
        10/16/2進表記、MSB/LSB、行/列優先、0/1反転など。
  \item \textbf{ピクセルエディタ}: 生成結果を1ドット単位で手動編集可能。
  \item \textbf{設定の自動保存と共有}: LocalStorage へ自動保存、
        URL クエリパラメータによる共有リンク生成。
\end{itemize}

\subsection{動作環境}

\begin{tabularx}{\textwidth}{l X}
\toprule
\textbf{項目} & \textbf{要件} \\
\midrule
ブラウザ & Chrome / Edge / Firefox / Safari の最新版 \\
JavaScript & 有効 \\
Canvas API & \texttt{CanvasRenderingContext2D}, \texttt{getImageData} 対応 \\
FontFace API & カスタムフォント読込に使用（主要ブラウザで対応済み）\\
LocalStorage & 設定の永続化に使用（オプトアウト可）\\
\bottomrule
\end{tabularx}

\newpage

\section{背景と課題}

\subsection{旧Python/Tkinter版の課題}

旧版（\texttt{font\_to\_binary\_GUI.py} / \texttt{font\_to\_binary\_CUI.py}）は
Python + Tkinter で実装された小規模ユーティリティであった。
一定の機能は果たしていたが、以下の課題があった。

\begin{itemize}[leftmargin=*,itemsep=2pt]
  \item \textbf{配布の煩雑さ}: Python と関連ライブラリ（Pillow, tkinter）を
        利用者がインストールする必要があった。
  \item \textbf{UI/UXの制約}: Tkinter のネイティブウィジェットでは
        近年の Web アプリのような細やかな表現ができず、
        レイアウトやテーマ切替も貧弱であった。
  \item \textbf{出力形式の限定}: C配列と一部の表記のみ対応しており、
        言語切替、ビットパック、MSB/LSB、0/1反転などの
        細かな設定ができなかった。
  \item \textbf{ピクセル編集不可}: 生成後の手動補正ができず、
        細かな調整をしたい場合は結果をコピーして手で書き換える必要があった。
  \item \textbf{共有性}: 設定を他人に渡す手段がなく、
        再現実験や共同作業に不向きであった。
\end{itemize}

\subsection{組み込み現場での典型的な需要}

マイコン用ディスプレイ（OLED/LCD/LED マトリクス）への
ドットフォント描画では、以下が繰り返し要求される。

\begin{itemize}[leftmargin=*,itemsep=2pt]
  \item ターゲットマイコン・表示ドライバごとに異なる
        バイト並び（行方向/列方向、MSB/LSB）への対応。
  \item FLASH 節約のためのビットパック配列出力（1bit/pixel）。
  \item \texttt{PROGMEM}, \texttt{const}, \texttt{static} 等の
        言語・処理系依存修飾子の付与。
  \item デザインを確認するためのプレビューと、
        生成結果の部分修正。
\end{itemize}

これらを毎回手作業で行うのは非効率であり、
\textbf{「設定をGUIで触り、結果をコピーするだけで済むツール」}
が常に求められていた。

\subsection{Webアプリ化の動機}

これらの課題に対し、
\textbf{クロスプラットフォームで即起動でき、機能拡張が容易で、
URLで設定を共有できる}という要求を満たすため、
Web アプリケーションとしての再実装を選択した。

\begin{itemize}[leftmargin=*,itemsep=2pt]
  \item 利用者はブラウザを開くだけで利用可能（インストール不要）。
  \item GitHub Pages による無料ホスティングで配布コスト \texttt{0}。
  \item React + TypeScript により型安全に機能拡張が可能。
  \item Canvas API + FontFace API によりフォントラスタライズを
        サーバレスで実行可能。
\end{itemize}

\newpage

\section{目的}

本ツールの目的は次の3点に集約される。

\begin{enumerate}[leftmargin=*,itemsep=2pt]
  \item \textbf{即時性}: ブラウザを開いてから配列を取得するまでの
        手数を最小化する。インストール、ログイン、ファイル入出力を要しない。
  \item \textbf{網羅性}: 組み込みで実際に遭遇する出力形式
        （言語・構造・ビット順・エンディアン・修飾子）を広くカバーする。
  \item \textbf{調整可能性}: 自動生成された結果を
        1ドット単位で手動編集できる手段を提供する。
\end{enumerate}

\subsection{達成基準}

\begin{tabularx}{\textwidth}{l X}
\toprule
\textbf{観点} & \textbf{基準} \\
\midrule
起動性 & ブラウザでURLを開くだけで利用可能、認証不要 \\
サイズ & 初期バンドル \texttt{< 500\,KB} gzipped（目標） \\
応答性 & 文字数100以下で描画再計算が \texttt{< 100\,ms}（目標） \\
対応文字 & Unicode 基本多言語面（BMP）全域、16×16 の想定でのみ品質保証 \\
多言語出力 & C/C++/Arduino/Python/JS/TS/Rust/Go/JSON/MD/Plain/Symbols の12形式 \\
出力構造 & 2次元/3次元/フラット/ビットパック（行・列）の5種類 \\
編集機能 & ペン/消しゴム/反転/回転/クリア/トグル \\
永続化 & LocalStorage 自動保存、URL 共有可能、オプトアウト可 \\
\bottomrule
\end{tabularx}

\subsection{非目標}

以下は本ツールのスコープ外とする。

\begin{itemize}[leftmargin=*,itemsep=2pt]
  \item アンチエイリアス付きのグレースケール出力
        （本ツールは2値ラスタライズに特化）。
  \item SVG や PNG などの画像ファイル形式での直接出力
        （必要ならプレビューを外部でキャプチャする）。
  \item ベクトルフォントの編集・グリフ追加機能。
  \item サーバーサイドでのフォント処理
        （すべてクライアントサイドで完結）。
  \item 商用フォントの再配布
        （同梱フォントはライセンスが明確なもののみ）。
\end{itemize}

\newpage

\section{アーキテクチャ}

\subsection{全体構成}

本ツールは純粋なクライアントサイド Single Page Application (SPA) であり、
サーバー通信は行わない（フォントファイル取得を除く）。

\begin{center}
\begin{tikzpicture}[
  node distance=3mm and 4mm,
  box/.style={draw=accentblue,fill=accentblue!8,rounded corners=1.5mm,
              minimum width=34mm,minimum height=10mm,align=center,
              font=\small},
  core/.style={draw=headerblue,fill=headerblue!10,rounded corners=1.5mm,
               minimum width=34mm,minimum height=10mm,align=center,
               font=\small\bfseries},
  arr/.style={-{Latex[length=2mm]},draw=headerblue,thick},
]
  \node[box] (user)   {ユーザー操作\\(UI)};
  \node[box,right=of user] (state)  {useAppState\\(useReducer)};
  \node[core,right=of state] (core) {core層\\(rasterize/encode/format)};
  \node[box,below=of user]  (persist) {usePersist\\(LocalStorage)};
  \node[box,below=of state] (share)   {share.ts\\(URL クエリ)};
  \node[box,below=of core]  (canvas)  {Canvas API\\FontFace API};
  \node[box,right=of core]  (output)  {OutputPanel\\(クリップボード)};

  \draw[arr] (user) -- (state);
  \draw[arr] (state) -- (core);
  \draw[arr] (core) -- (output);
  \draw[arr] (state) -- (persist);
  \draw[arr] (state) -- (share);
  \draw[arr] (core) -- (canvas);
\end{tikzpicture}
\end{center}

\subsection{レイヤ分割}

\begin{tabularx}{\textwidth}{l X}
\toprule
\textbf{レイヤ} & \textbf{責務} \\
\midrule
\texttt{src/core/} & 純粋関数群。DOM依存を最小化（Canvas利用箇所のみ）。ラスタライズ、ビットパック、フォーマッタ、共有URLエンコード。\\
\texttt{src/hooks/} & React フック。状態管理（useReducer）、永続化（LocalStorage）、テーマ切替、グリフメモ化。\\
\texttt{src/components/} & UI コンポーネント。shadcn/ui プリミティブを基盤とし、パネル単位で責務分離。\\
\texttt{src/components/ui/} & Radix UI ベースの汎用プリミティブ（Button, Card, Input, Tabs, Switch, Slider, Select, Tooltip 等）。\\
\bottomrule
\end{tabularx}

\subsection{データフロー}

\begin{enumerate}[leftmargin=*,itemsep=2pt]
  \item ユーザー入力 → \texttt{useAppState} の dispatch でアクション発行。
  \item Reducer が新しい \texttt{AppState} を返し、コンポーネントが再描画。
  \item \texttt{useGlyphs} が入力テキストとフォント設定に基づき
        \texttt{rasterizeChar()} を再実行しグリフ配列をメモ化。
  \item \texttt{formatOutput()} がグリフ配列 + フォーマット設定から
        出力文字列を生成し、\texttt{OutputPanel} が表示。
  \item 同時に \texttt{usePersist} が \texttt{AppState} を
        LocalStorage に書き戻す（\texttt{overrides} を除く）。
\end{enumerate}

\subsection{ビルド・配布パイプライン}

\begin{enumerate}[leftmargin=*,itemsep=2pt]
  \item 開発: \texttt{npm run dev}（Vite 開発サーバ、HMR有効）。
  \item 本番ビルド: \texttt{npm run build}（Vite が \texttt{dist/} を生成）。
  \item 自動デプロイ: \texttt{main} ブランチ push を契機に
        \texttt{.github/workflows/deploy.yml} が起動、
        GitHub Pages（\texttt{gh-pages} 相当の Actions 経由）へ配信。
  \item SPA フォールバック: デプロイ時に \texttt{404.html} を
        \texttt{index.html} と同内容でコピーし、
        任意パスでのアクセスに対応。
\end{enumerate}

\newpage

\section{ファイル構成}

\subsection{ディレクトリツリー}

\begin{lstlisting}[basicstyle=\ttfamily\footnotesize,language={},numbers=none]
font_to_bin/
├── public/
│   ├── fonts/
│   │   ├── DotGothic16-Regular.ttf
│   │   └── misaki_gothic_2nd.ttf
│   ├── favicon.svg
│   └── manifest.webmanifest
├── src/
│   ├── core/
│   │   ├── types.ts       … 型定義（AppState, Glyph, Format 他）
│   │   ├── fonts.ts       … 同梱フォント定義 + カスタム登録API
│   │   ├── rasterize.ts   … Canvas によるグリフ2値化
│   │   ├── encode.ts      … ビットパック・数値フォーマット
│   │   ├── format.ts      … 言語別出力文字列生成
│   │   ├── share.ts       … URLクエリパラメータ共有
│   │   ├── presets.ts     … 用途別プリセット(FormatOptions雛形)
│   │   └── defaults.ts    … 既定値
│   ├── components/
│   │   ├── ui/            … shadcn/ui プリミティブ群
│   │   ├── Header.tsx
│   │   ├── EasyPanel.tsx        … かんたんモード(3ステップ)
│   │   ├── InputPanel.tsx
│   │   ├── FormatPanel.tsx
│   │   ├── PreviewPanel.tsx     … 横スクロール帯プレビュー
│   │   ├── GlyphPreview.tsx     … 補助線対応のグリフ描画
│   │   ├── OutputPanel.tsx
│   │   ├── PixelEditor.tsx
│   │   └── FreeCanvasPanel.tsx  … 自由モードのキャンバス
│   ├── hooks/
│   │   ├── useAppState.ts … useReducer ラッパ
│   │   ├── useGlyphs.ts   … メモ化ラスタライザ
│   │   ├── usePersist.ts  … LocalStorage自動保存
│   │   ├── useMode.ts     … かんたん/詳細モード切替
│   │   └── useTheme.ts    … ダーク/ライト
│   ├── lib/utils.ts       … cn() ヘルパ
│   ├── App.tsx            … 3カラムレイアウト + モーダル制御
│   ├── main.tsx           … エントリポイント
│   └── index.css          … Tailwind ベースCSS
├── .github/workflows/deploy.yml   … GitHub Pages 自動デプロイ
├── doc/                           … LaTeX 設計仕様書
├── legacy/                        … 旧Python版（参考保存）
├── index.html
├── package.json
├── tsconfig.json
├── tsconfig.node.json
├── vite.config.ts
├── tailwind.config.js
├── postcss.config.js
└── README.md
\end{lstlisting}

\subsection{core 層}

\begin{itemize}[leftmargin=*,itemsep=2pt]
  \item \texttt{src/core/types.ts}\,---\,\texttt{AppState}, \texttt{Glyph},
        \texttt{FormatOptions} など全型定義。UI と core の契約。
  \item \texttt{src/core/fonts.ts}\,---\,\texttt{BUILTIN\_FONTS} 配列、
        \texttt{loadFont}, \texttt{registerCustomFont}。
        FontFace API でカスタム TTF/OTF を登録。
  \item \texttt{src/core/rasterize.ts}\,---\,\texttt{rasterizeChar(char, opts)}
        $\rightarrow$ \texttt{Glyph}。Canvas へ1文字描画し
        \texttt{getImageData} でアルファ値を閾値2値化。
  \item \texttt{src/core/encode.ts}\,---\,\texttt{packRow} / \texttt{packColumn}
        （MSB/LSB ビットパック）、\texttt{formatNumber}（bin/hex/dec）。
  \item \texttt{src/core/format.ts}\,---\,\texttt{formatOutput}。
        出力言語に応じて \texttt{formatC} / \texttt{formatPython}
        / \texttt{formatJS} / \texttt{formatRust} / \texttt{formatGo}
        / \texttt{formatJSON} / \texttt{formatMarkdown}
        / \texttt{formatSymbols} / \texttt{formatPlain} を呼び分け。
  \item \texttt{src/core/share.ts}\,---\,\texttt{encodeStateToURL}
        / \texttt{decodeStateFromURL}。
        \texttt{AppState} を base64url 化して \texttt{?s=...} に載せる。
  \item \texttt{src/core/presets.ts}\,---\,\texttt{PRESETS} 配列
        （Arduino/C言語, C言語2D, Python リスト, MicroPython bytes,
        JSON, 記号アートの6種）と \texttt{detectPreset()}
        （現在の \texttt{FormatOptions} から該当プリセットIDを推定）。
        かんたんモードでワンタッチ切替に使用。
  \item \texttt{src/core/defaults.ts}\,---\,\texttt{DEFAULT\_STATE}
        （text=\texttt{"ABC"}, fontId=\texttt{dotgothic16},
        16$\times$16, threshold=128）。
\end{itemize}

\subsection{hooks 層}

\begin{itemize}[leftmargin=*,itemsep=2pt]
  \item \texttt{src/hooks/useAppState.ts}\,---\,\texttt{useReducer}。
        Action は \texttt{setText}, \texttt{setSize},
        \texttt{setFormat}, \texttt{setOverride} のほか、
        プレビュー補助線用 \texttt{setPreviewGridStep}、
        自由モード用 \texttt{setFreeSize} / \texttt{setFreeMatrix} /
        \texttt{setFreePixel} / \texttt{setFreeGridStep} /
        \texttt{clearFree} を提供する。
  \item \texttt{src/hooks/useGlyphs.ts}\,---\,text / font / size
        変更時のみ再ラスタライズするメモ化。overrides をマージ。
  \item \texttt{src/hooks/usePersist.ts}\,---\,LocalStorage キー
        \texttt{font-to-bin.state.v1}。
  \item \texttt{src/hooks/useMode.ts}\,---\,かんたん/詳細/自由モード切替。
        LocalStorage キー \texttt{font\_to\_bin\_mode\_v1}。値は
        \texttt{easy} / \texttt{advanced} / \texttt{free} のいずれか。
  \item \texttt{src/hooks/useTheme.ts}\,---\,light / dark 切替、
        \texttt{font-to-bin.theme} キー。
\end{itemize}

\subsection{UI 層}

\begin{itemize}[leftmargin=*,itemsep=2pt]
  \item \texttt{src/components/Header.tsx}\,---\,タイトル、
        \textbf{かんたん/詳細モード切替}のセグメントタブ、
        テーマ切替、URL 共有、GitHub リンク。
  \item \texttt{src/components/EasyPanel.tsx}\,---\,
        かんたんモード用の3ステップUI。
        \textbf{ステップ1}: 変換したい文字。
        \textbf{ステップ2}: フォント + ドットサイズクイックボタン
        (8/12/16/24/32px) + 太字トグル。
        \textbf{ステップ3}: 用途プリセットを6種から1クリック選択
        （\texttt{presets.ts} の \texttt{PRESETS} を使用）。
  \item \texttt{src/components/InputPanel.tsx}\,---\,詳細モードの左上。
        テキスト、フォント選択、幅/高さ、描画オフセット、太字、閾値。
  \item \texttt{src/components/FormatPanel.tsx}\,---\,出力言語、
        構造、基数、MSB/LSB、0/1 反転、変数名、データ型。
  \item \texttt{src/components/PreviewPanel.tsx}\,---\,グリフ一覧を
        \textbf{横スクロール帯}で表示。補助線間隔セレクタ、
        ズーム（50--400\%）を備える。
        文字クリックでピクセルエディタを開く。
  \item \texttt{src/components/GlyphPreview.tsx}\,---\,単一グリフの
        Canvas タイル描画。\texttt{gridStep} プロパティで N ドット毎の
        強調補助線を描画。
  \item \texttt{src/components/OutputPanel.tsx}\,---\,生成結果の
        表示 + コピー + ダウンロード。
  \item \texttt{src/components/PixelEditor.tsx}\,---\,モーダル。
        ペン/消しゴム/トグル/反転/回転/クリア/ズームに加え、
        補助線間隔の切替に対応。
  \item \texttt{src/components/FreeCanvasPanel.tsx}\,---\,自由モード専用。
        W$\times$H(1--256)指定、スクロール可能なキャンバス、補助線、
        ズーム、ペン／消しゴム／トグル、反転・回転・全消去。
        編集結果は \texttt{AppState.free.matrix} に
        \texttt{setFreePixel}（1セル単位）で反映され、
        出力側では単一グリフとして扱う。
  \item \texttt{src/components/ui/}\,---\,Radix UI ベースの
        shadcn/ui 汎用プリミティブ群。
\end{itemize}

\subsection{その他}

\begin{itemize}[leftmargin=*,itemsep=2pt]
  \item \texttt{.github/workflows/deploy.yml}\,---\,Node 20 で
        \texttt{npm ci} $\rightarrow$ \texttt{build} $\rightarrow$
        Pages にアップロード。
  \item \texttt{public/fonts/*.ttf}\,---\,同梱フォント
        （DotGothic16, Misaki Gothic 2nd）。
  \item \texttt{legacy/}\,---\,旧 Python/Tkinter 実装
        （参照用、ビルドからは除外）。
\end{itemize}

\newpage

\section{技術スタック}

\subsection{選定方針}

\begin{itemize}[leftmargin=*,itemsep=2pt]
  \item \textbf{型安全と開発体験}: TypeScript + Vite による高速な
        HMR と厳格な型チェック。
  \item \textbf{標準準拠}: Canvas 2D / FontFace / LocalStorage など
        W3C 標準 API のみを使用し、フレームワーク依存を最小化。
  \item \textbf{UI 構築速度}: Tailwind CSS + shadcn/ui（Radix UI ベース）
        でスタイルと挙動を分離しつつ、アクセシビリティを確保。
  \item \textbf{ゼロ運用コスト}: GitHub Pages による静的配信。
        Actions で自動デプロイ、サーバは持たない。
\end{itemize}

\subsection{使用ライブラリ・ツール}

\begin{tabularx}{\textwidth}{l l X}
\toprule
\textbf{分類} & \textbf{名称} & \textbf{用途} \\
\midrule
言語       & TypeScript 5     & 型安全な JavaScript。\\
ランタイム & Node.js 20+      & 開発・ビルド環境。\\
ビルド     & Vite 6           & 開発サーバ + 本番バンドル。\\
UI 基盤    & React 18         & 宣言的 UI、Hooks ベース状態管理。\\
スタイル   & Tailwind CSS 3   & ユーティリティファースト CSS。\\
UI 部品    & shadcn/ui        & Radix UI + Tailwind プリミティブ集。\\
UI プリミティブ & Radix UI    & アクセシブルな Tabs/Select/Slider 等。\\
アイコン   & lucide-react     & SVG アイコン。\\
CVA        & class-variance-authority & バリアント管理。\\
\bottomrule
\end{tabularx}

\subsection{ブラウザ API}

\begin{tabularx}{\textwidth}{l X}
\toprule
\textbf{API} & \textbf{役割} \\
\midrule
Canvas 2D (\texttt{getContext("2d")})     & フォント描画とピクセル抽出。\\
\texttt{getImageData}                     & RGBA 配列取得、アルファ閾値 2 値化の基礎。\\
\texttt{FontFace}                         & カスタム TTF/OTF/WOFF の動的登録。\\
\texttt{URL} / \texttt{URLSearchParams}   & 共有リンクのクエリパラメータ読み書き。\\
\texttt{LocalStorage}                     & 設定の永続化。\\
\texttt{navigator.clipboard}              & 出力結果のコピー。\\
\texttt{matchMedia("prefers-color-scheme")} & OS テーマ検出。\\
\bottomrule
\end{tabularx}

\subsection{同梱フォント}

\begin{tabularx}{\textwidth}{l l X}
\toprule
\textbf{フォント} & \textbf{推奨サイズ} & \textbf{ライセンス} \\
\midrule
DotGothic16       & 16$\times$16 & SIL OFL 1.1 \\
Misaki Gothic 2nd & 8$\times$8   & 自由配布（商用含む）\\
システムモノスペース & 任意       & ブラウザ標準フォント依存 \\
\bottomrule
\end{tabularx}

\subsection{ビルド設定}

\begin{itemize}[leftmargin=*,itemsep=2pt]
  \item \texttt{vite.config.ts}: \texttt{base} を本番時のみ
        \texttt{/font\_to\_bin/} に切り替え。
  \item \texttt{tsconfig.json}: \texttt{strict: true}, \texttt{noUnusedLocals},
        \texttt{noUnusedParameters} を有効化。
  \item \texttt{tailwind.config.js}: \texttt{darkMode: "class"}、
        カスタムカラートークンを定義。
  \item PWA manifest (\texttt{public/manifest.webmanifest}):
        アイコン・テーマカラー・表示モードを定義（SW は持たない）。
\end{itemize}

\newpage

\section{API 仕様}

本章では \texttt{src/core/} が提供する公開 API と主要な型を示す。
UI 層はこれらの関数・型にのみ依存し、DOM 操作は core 内の
\texttt{rasterize.ts} に閉じ込められている。

\subsection{型定義（types.ts）}

\begin{apibox}{AppState}
\texttt{AppState} は UI 全体の単一情報源であり、永続化・共有の単位でもある。

\begin{lstlisting}
type AppState = {
  text: string;
  fontId: string;
  width: number;        // グリフ幅[px]
  height: number;       // グリフ高さ[px]
  size: number;         // 描画フォントサイズ[px]
  threshold: number;    // 2値化閾値 0-255
  offsetX: number;
  offsetY: number;
  bold: boolean;
  format: FormatOptions;
  overrides: Record<number, Matrix>; // charIndex→手動編集結果
};
\end{lstlisting}
\end{apibox}

\begin{apibox}{Glyph / Matrix}
\begin{lstlisting}
type Bit = 0 | 1;
type Matrix = Bit[][]; // matrix[row][col]
type Glyph = {
  char: string;
  width: number;
  height: number;
  matrix: Matrix;
};
\end{lstlisting}
\end{apibox}

\begin{apibox}{FormatOptions}
\begin{lstlisting}
type FormatOptions = {
  language: OutputLanguage;  // c|cpp|arduino|python|...
  structure: Structure;      // matrix2d|matrix3d|flat|bitpack-row|bitpack-col
  dataType: string;          // "uint8_t" 等
  bitOrder: BitOrder;        // "msb" | "lsb"
  radix: Radix;              // "bin" | "hex" | "dec"
  invert: boolean;
  variableName: string;
  itemsPerLine: number;
  includeCharComment: boolean;
  includeAsciiArt: boolean;
  padByte: boolean;
  indent: string;
  prefix: string;
  suffix: string;
};
\end{lstlisting}
\end{apibox}

\subsection{ラスタライズ（rasterize.ts）}

\begin{funcblock}{rasterizeChar(char: string, opts: RasterizeOptions): Glyph}
指定文字を Canvas に描画し、各ピクセルのアルファ値を
\texttt{opts.threshold} で 2 値化して \texttt{Glyph} を返す。

\apisection{引数}
\begin{itemize}[leftmargin=*,itemsep=1pt]
  \item \texttt{char}: 対象文字（1文字）。
  \item \texttt{opts.fontFamily}: CSS の \texttt{font-family}。
  \item \texttt{opts.width}, \texttt{opts.height}: 出力行列のサイズ。
  \item \texttt{opts.size}: 描画フォントの \texttt{px} サイズ。
  \item \texttt{opts.threshold}: $[0,255]$ のアルファ閾値。
  \item \texttt{opts.offsetX}, \texttt{opts.offsetY}: 描画オフセット。
  \item \texttt{opts.bold}: \texttt{true} で \texttt{bold} 書体。
\end{itemize}

\apisection{戻り値}
\texttt{Glyph}。\texttt{matrix[row][col]} は 0/1 の \texttt{Bit}。
\end{funcblock}

\subsection{エンコード（encode.ts）}

\begin{funcblock}{packRow(matrix: Matrix, bitOrder: BitOrder, padByte: boolean): number[]}
\texttt{matrix} を行方向に走査して 1bit/pixel にパックする。
\texttt{padByte=true} の場合、各行を 8bit 境界にゼロ埋めする。
\end{funcblock}

\begin{funcblock}{packColumn(matrix: Matrix, bitOrder: BitOrder, padByte: boolean): number[]}
列方向パック版。OLED/LCD の「縦8ドット=1バイト」フォーマットに相当。
\end{funcblock}

\begin{funcblock}{formatNumber(n: number, radix: Radix, bitWidth: number): string}
\texttt{n} を指定基数で文字列化する。\texttt{hex} は \texttt{0x} 接頭辞、
\texttt{bin} は \texttt{0b} 接頭辞、\texttt{dec} は素の 10 進。
\texttt{bitWidth} に応じて 0 埋めする。
\end{funcblock}

\subsection{フォーマット（format.ts）}

\begin{funcblock}{formatOutput(glyphs: Glyph[], opts: FormatOptions): string}
グリフ列とフォーマット設定から、ユーザーに表示する最終的な
出力文字列を生成する。言語別に以下のヘルパへディスパッチする。

\apisection{ディスパッチ先}
\begin{itemize}[leftmargin=*,itemsep=1pt]
  \item \texttt{formatC} (C / C++ / Arduino)
  \item \texttt{formatPython}
  \item \texttt{formatJS} (JavaScript / TypeScript)
  \item \texttt{formatRust}
  \item \texttt{formatGo}
  \item \texttt{formatJSON}
  \item \texttt{formatMarkdown}
  \item \texttt{formatSymbols} (記号アート、\texttt{■/□})
  \item \texttt{formatPlain} (生の 0/1 テキスト)
\end{itemize}
\end{funcblock}

\subsection{フォント管理（fonts.ts）}

\begin{funcblock}{loadFont(font: FontDef): Promise<void>}
\texttt{FontDef.url} が指す TTF/OTF を \texttt{FontFace} で読み込み、
\texttt{document.fonts} に登録する。同梱フォント用。
\end{funcblock}

\begin{funcblock}{registerCustomFont(file: File): Promise<FontDef>}
ユーザーがアップロードした \texttt{File} を \texttt{FontFace} 経由で
登録し、ID を付与した \texttt{FontDef} を返す。
以降の \texttt{fontId} 指定で利用可能になる。
\end{funcblock}

\subsection{共有リンク（share.ts）}

\begin{funcblock}{encodeStateToURL(state: AppState): string}
\texttt{overrides} を除く \texttt{AppState} を JSON 化し、base64url で
圧縮した文字列を \texttt{?s=} クエリに載せた URL を返す。
\end{funcblock}

\begin{funcblock}{decodeStateFromURL(): Partial<AppState> | null}
現在の URL から \texttt{?s=} を読み出し、復元した \texttt{AppState} を返す。
不正な値や欠損フィールドは無視し、欠けた箇所は既定値で補完される。
\end{funcblock}

\begin{notebox}[overrides を共有しない理由]
ピクセル編集結果（\texttt{overrides}）はサイズが大きく URL に載せづらいため、
共有リンクには含めない。共有先で再生成するか、
JSON エクスポート機能（\texttt{OutputPanel} のダウンロード）を用いる。
\end{notebox}

\subsection{React Hooks}

\begin{funcblock}{useAppState(initial: AppState): [AppState, Dispatch<Action>]}
\texttt{useReducer} を用いた状態管理。サポートするアクション:
\texttt{setText}, \texttt{setFontId}, \texttt{setSize},
\texttt{setWidth}, \texttt{setHeight}, \texttt{setFontSize},
\texttt{setThreshold}, \texttt{setOffsetX}, \texttt{setOffsetY},
\texttt{setBold}, \texttt{setFormat}, \texttt{setOverride},
\texttt{clearOverride}, \texttt{clearAllOverrides}, \texttt{replaceState}。
\end{funcblock}

\begin{funcblock}{useGlyphs(state: AppState): Glyph[]}
\texttt{text}, \texttt{fontId}, 幾何パラメータ、\texttt{overrides}
の変化にのみ反応して再計算するメモ化ラッパ。
\end{funcblock}

\begin{funcblock}{usePersist(state: AppState)}
\texttt{AppState} を \texttt{localStorage["font-to-bin.state.v1"]} に
自動保存する副作用フック。\texttt{overrides} は除外する。
\end{funcblock}

\newpage

\section{実行モデル}

\subsection{起動シーケンス}

\begin{enumerate}[leftmargin=*,itemsep=2pt]
  \item \texttt{index.html} が \texttt{src/main.tsx} を読み込み、
        React ルートを \texttt{\#root} にマウント。
  \item \texttt{App.tsx} が初期 \texttt{AppState} を決定する。
    \begin{itemize}[leftmargin=*,itemsep=1pt]
      \item URL に \texttt{?s=...} があれば \texttt{decodeStateFromURL()} を
            使用して復元。
      \item なければ LocalStorage から \texttt{persistedInitial()} で読み込み。
      \item どちらも無効な場合は \texttt{DEFAULT\_STATE} を採用。
    \end{itemize}
  \item \texttt{useGlyphs} が初期テキストをラスタライズし、
        \texttt{formatOutput} の結果が \texttt{OutputPanel} に反映される。
  \item \texttt{usePersist} が以後の \texttt{AppState} 変化を
        LocalStorage に書き込み続ける。
\end{enumerate}

\subsection{ラスタライズの詳細}

\begin{enumerate}[leftmargin=*,itemsep=2pt]
  \item \texttt{<canvas>} を $\texttt{width}\times\texttt{height}$ で作成。
  \item 背景を白でクリアし、\texttt{fillStyle="black"} で文字を描画。
    \begin{itemize}[leftmargin=*,itemsep=1pt]
      \item \texttt{ctx.font = `\${size}px \${fontFamily}`}
            （太字時は \texttt{bold} を前置）。
      \item \texttt{ctx.textBaseline = "top"}。
      \item \texttt{fillText(char, offsetX, offsetY)}。
    \end{itemize}
  \item \texttt{getImageData(0,0,width,height)} で RGBA 配列を取得。
  \item 各ピクセルのアルファチャネル $\alpha$ と閾値 $\tau$ を比較し、
        $\alpha \ge \tau$ を $1$、それ以外を $0$ とする。
  \item 結果を \texttt{Matrix} （\texttt{matrix[row][col]}）に格納。
\end{enumerate}

\subsection{ビットパック規則}

\subsubsection{行方向パック（bitpack-row）}

\begin{itemize}[leftmargin=*,itemsep=1pt]
  \item 各行を左$\to$右に走査し、8bit 単位で 1 バイトを構成。
  \item \texttt{bitOrder="msb"}: 左端ピクセルが bit7。
  \item \texttt{bitOrder="lsb"}: 左端ピクセルが bit0。
  \item \texttt{padByte=true} 時、行終端が 8bit 境界でなければ
        下位（または上位）を 0 埋め。
\end{itemize}

\subsubsection{列方向パック（bitpack-col）}

\begin{itemize}[leftmargin=*,itemsep=1pt]
  \item 各列を上$\to$下に走査し、縦 8bit を 1 バイトに格納。
        OLED SSD1306 の \texttt{horizontal addressing mode} と整合。
  \item \texttt{bitOrder="msb"}: 上端ピクセルが bit7。
  \item \texttt{bitOrder="lsb"}: 上端ピクセルが bit0（SSD1306 既定）。
\end{itemize}

\subsection{状態遷移（Reducer）}

\texttt{useAppState} の Reducer は次の不変条件を守る。

\begin{itemize}[leftmargin=*,itemsep=2pt]
  \item \texttt{width} / \texttt{height} / \texttt{size} の変更時、
        既存の \texttt{overrides} は次元が一致しないため \texttt{clearAllOverrides} 扱い。
  \item \texttt{text} の変更時、文字数が減少すれば末尾の \texttt{overrides} を破棄。
  \item \texttt{fontId} の変更時、推奨サイズがあれば \texttt{size} を合わせる
        （ユーザーが明示的に変更中でない場合のみ）。
  \item \texttt{replaceState}（共有リンク適用や読込）は \texttt{overrides} を空にする。
\end{itemize}

\subsection{永続化}

\begin{itemize}[leftmargin=*,itemsep=2pt]
  \item キー: \texttt{font-to-bin.state.v1}（スキーマ変更時は \texttt{.v2} へ上げる）。
  \item 保存対象: \texttt{AppState} から \texttt{overrides} を除いた部分。
  \item 復元時のバリデーション: 型不一致・未知フィールドは無視し、
        欠損は既定値で補完。
  \item ユーザーはヘッダーの「設定リセット」で LocalStorage を削除可能。
\end{itemize}

\subsection{共有リンクのラウンドトリップ}

\begin{enumerate}[leftmargin=*,itemsep=2pt]
  \item 「URL共有」ボタン押下で \texttt{encodeStateToURL(state)} を実行。
  \item \texttt{JSON.stringify} $\rightarrow$ \texttt{TextEncoder}
        $\rightarrow$ base64url 変換。
  \item 現在の URL の \texttt{?s=} を差し替え、クリップボードへコピー。
  \item 開いた側では起動時に \texttt{decodeStateFromURL()} が復元。
        展開失敗時は既定値へフォールバック。
\end{enumerate}

\subsection{パフォーマンス考慮}

\begin{itemize}[leftmargin=*,itemsep=2pt]
  \item \texttt{useGlyphs} は \texttt{text + geometry + overrides}
        のハッシュキー変化時のみ再実行。
  \item フォーマッタは文字数 $N$、幅 $W$、高さ $H$ に対して $O(NWH)$。
        典型値（$N=8, W=H=16$）では数百~$\mu$s。
  \item ピクセルエディタは 1 グリフあたり $W\times H$ ノードで描画するが、
        SVG の仮想 DOM に任せており数ミリ秒以内で再描画可能。
\end{itemize}

\subsection{デプロイ}

\begin{itemize}[leftmargin=*,itemsep=2pt]
  \item \texttt{main} ブランチへの push で GitHub Actions が起動。
  \item \texttt{actions/setup-node@v4} $\rightarrow$ \texttt{npm ci}
        $\rightarrow$ \texttt{npm run build} $\rightarrow$
        \texttt{actions/upload-pages-artifact}
        $\rightarrow$ \texttt{actions/deploy-pages}。
  \item ビルド後、\texttt{dist/index.html} を \texttt{404.html} にコピーして
        SPA の任意パス直打ちに対応。
\end{itemize}

\begin{cautionbox}[初回のみ必要な設定]
リポジトリの \textbf{Settings $\to$ Pages $\to$ Source} を
\texttt{GitHub Actions} に設定する必要がある（\texttt{gh-pages} ブランチは使わない）。
\end{cautionbox}



\end{document}
